\section*{A. Introduction}
In this paper we continue the study of Hodge ideals initiated in \cite{MP1}, \cite{MP2}, by considering an analogous theory 
for arbitrary $\jepPubBold{Q}$-divisors. The emphasis here is on a birational definition and study of Hodge ideals, while the companion 
paper \cite{MP3} is devoted to a study based on their connection with the $V$-filtration, inspired by \cite{Saito-MLCT}. Both 
approaches turn out to provide crucial information towards a complete understanding of these objects.

Let $X$ be a smooth complex variety. If $D$ is reduced divisor on $X$, the Hodge ideals $I_k (D)$, with $k \geqslant 0$, 
are defined in terms of the Hodge filtration 
on the $\jepPubScript{D}_X$-module $\jepPubScript{O}_X(*D)$ of functions with poles of arbitrary order along $D$. Indeed, this 
$\jepPubScript{D}_X$-module
underlies a mixed Hodge module on $X$, and therefore comes with a Hodge filtration $F_\jepPubBullet \jepPubScript{O}_X (*D)$, which satisfies
$$F_k \jepPubScript{O}_X(*D) = I_k (D) \otimes \jepPubScript{O}_X\big( (k+1)D\big)\quad\text{for all}\,\,\,\, k \geqslant 0.$$
See \cite{MP1} for details, and for an extensive study of the ideals $I_k (D)$.

Our goal here is to provide a similar construction and study in the general case. A natural device for dealing with the 
fact that fractional divisors are not directly related to Hodge theory is to use new objects derived from covering constructions.
Let $D$ be an arbitrary effective $\jepPubBold{Q}$-divisor on $X$. Locally, we can write $D = \alpha H$, for 
some $\alpha \in \jepPubBold{Q}_{> 0}$ and $H = {\rm div} (h)$, the divisor of a nonzero regular function; we also denote by $Z$ the support of $D$.
A well-known construction associates to this data a twisted version of the localization $\jepPubScript{D}$-module above, namely
$$\jepPubScript{M} (h^{-\alpha}) : = \jepPubScript{O}_X (*Z) h^{-\alpha},$$
that is the rank $1$ free $\jepPubScript{O}_X(*Z)$-module 
with generator the symbol $h^{-\alpha}$, on which a derivation $D$ of $\jepPubScript{O}_X$ acts by
$$D(wh^{-\alpha}) :=\Big(D(w) - \alpha w\frac{ D(h)}{h}\Big)h^{-\alpha}.$$
It turns out that this $\jepPubScript{D}_X$-module can be endowed with a natural filtration $F_k \jepPubScript{M} (h^{-\alpha})$, with 
$k \geqslant 0$, which makes it a filtered direct summand of a $\jepPubScript{D}$-module underlying a mixed Hodge module on $X$; see 
Section~\ref{coverings}. This plays a role 
analogous to the Hodge filtration, and just as in the reduced case one can show that $F_k \jepPubScript{M} (h^{-\alpha}) \subseteq 
\jepPubScript{O}_X (kZ ) h^{-\alpha}$. This is done in Section~\ref{smooth} and Section~\ref{definition}, by first analyzing the case when $Z$ is a smooth divisor 
(in this case, if $\lceil D\rceil=Z$, then
the inclusion is in fact an equality). It is therefore natural to define the \emph{$k$-th Hodge ideal} of $D$ by the formula
$$F_k \jepPubScript{M} (h^{-\alpha}) = I_k (D) \otimes_{\jepPubScript{O}_X} \jepPubScript{O}_X (kZ ) h^{-\alpha}.$$ 

Similarly to \cite{MP1}, one of our main goals here is to study Hodge ideals of $\jepPubBold{Q}$-divisors by means of log resolutions. 
To this end, let $f\colon Y\to X$ be a log resolution of the pair $(X, D)$ that is an isomorphism over $U=X\smallsetminus Z$,  
and denote $g=h\circ f$. There is a filtered isomorphism
$$
\big(\jepPubScript{M}(h^{-\alpha}), F \big)\simeq f_+\big(\jepPubScript{M}(g^{-\alpha}), F\big).
$$
Denoting $G=f^*D$ and $E = {\rm Supp}(G)$, so that $E$ is a simple normal crossing divisor, 
it turns out that there exists a complex on $Y$:
$$C^{\jepPubBullet}_{g^{-\alpha}} (- \lceil G\rceil):\,\,0\longrightarrow\jepPubScript{O}_Y(-\lceil G\rceil)\otimes_{\jepPubScript{O}_Y} \jepPubScript{D}_Y\longrightarrow\jepPubScript{O}_Y(-\lceil G\rceil)\otimes_{\jepPubScript{O}_Y}\Omega^1_Y(\log E)\otimes_{\jepPubScript{O}_Y}\jepPubScript{D}_Y$$
$$\longrightarrow\cdots\longrightarrow\jepPubScript{O}_Y(-\lceil G\rceil)\otimes_{\jepPubScript{O}_Y}\omega_Y(E)\otimes_{\jepPubScript{O}_Y}\jepPubScript{D}_Y\longrightarrow 0,$$
which is placed in degrees $-n,\ldots,0$, whose differential is described in Section~\ref{complex_for_SNC}.
This complex has a natural filtration given, for $k\geqslant 0$, by subcomplexes 
$$F_{k-n}C^{\jepPubBullet}_{g^{-\alpha}} (- \lceil G\rceil): = 
0 \rightarrow \jepPubScript{O}_Y(-\lceil G\rceil)\otimes  F_{k-n} \jepPubScript{D}_Y$$
$$\rightarrow  \jepPubScript{O}_Y(-\lceil G\rceil)\otimes\Omega_Y^1(\log E) \otimes  F_{k-n+1} \jepPubScript{D}_Y \rightarrow 
\cdots \rightarrow  \jepPubScript{O}_Y(-\lceil G\rceil)\otimes \omega_Y(E) \otimes  F_k \jepPubScript{D}_Y\rightarrow 0.$$

\noindent
Extending \cite[Prop. 3.1]{MP1}, we show in Proposition \ref{prop_SNC_complex} and Proposition \ref{Hodge_ideals_SNC} 
that there is a filtered quasi-isomorphism
$$
\big(C^{\jepPubBullet}_{g^{-\alpha}} (- \lceil G\rceil), F \big) \simeq \big(\jepPubScript{M}_r (g^{-\alpha}), F \big),
$$
where $\jepPubScript{M}_r (g^{-\alpha})$ is the filtered right $\jepPubScript{D}_Y$-module associated to $\jepPubScript{M} (g^{-\alpha})$.
Thus one can use $\big(C^{\jepPubBullet}_{g^{-\alpha}} (- \lceil G\rceil), F\big)$ as a concrete representative for computing the filtered 
$\jepPubScript{D}$-module pushforward of $\big(\jepPubScript{M}_r (g^{-\alpha}), F\big)$, hence for computing 
the ideals $I_k (D)$. More precisely, we have 
$$R^0f_*F_{k-n}\big(C^{\jepPubBullet}_{g^{-\alpha}} (- \lceil G\rceil)\otimes_{\jepPubScript{D}_Y}\jepPubScript{D}_{Y\to X}\big)\simeq
h^{-\alpha}\omega_X (kZ )\otimes_{\jepPubScript{O}_X}I_k(D).$$
See Theorem \ref{formula_log_resolution} for a complete picture regarding this push-forward operation.

This fact, together with special properties of the filtration on $\jepPubScript{D}$-modules underlying mixed Hodge modules, leads to 
our main results on Hodge ideals, which are collected in the following:

\begin{intro-theorem}\label{main_list}
In the set-up above, the Hodge ideals $I_k (D)$ satisfy:


\noindent
(i) $I_0 (D)$ is the multiplier ideal $\jepPubScript{I} \big((1 - \varepsilon)D\big)$, 
so in particular $I_0 (D) = \jepPubScript{O}_X$ if and only if the pair $(X, D)$ is log canonical; see Section~9.


\noindent
(ii) If $Z$ has simple normal crossings, then 
$$I_k (D) = I_k (Z) \otimes \jepPubScript{O}_X (Z - \lceil D \rceil),$$
while $I_k (Z)$ can be computed explicitly as in \cite[Prop.~8.2]{MP1}; see Section~\ref{SNC}.  In particular, if 
$Z$ is smooth, then $I_k (D) = \jepPubScript{O}_X (Z - \lceil D \rceil)$ for all $k$; cf. also Corollary 11.12.



\noindent
(iii) The Hodge filtration is generated at level $n-1$, where $n = \dim X$, i.e. 
$$F_\ell \jepPubScript{D}_X\cdot\big(I_k(D)\otimes \jepPubScript{O}_X(kZ)h^{-\alpha}\big) = I_{k+\ell}(D)\otimes \jepPubScript{O}_X \big((k+\ell)Z\big)h^{-\alpha}$$
for all $k \geqslant n-1$ and $\ell\geqslant 0$; see Section~10.


\noindent
(iv) There are non-triviality criteria for $I_k (D)$ at a point $x \in D$ in terms of the multiplicity of $D$ at $x$; 
see Section~11.



\noindent
(v) If $X$ is projective, $I_k(D)$ satisfy a vanishing theorem analogous to Nadel Vanishing for multiplier ideals; 
see Section~12.


\noindent
(vi) If $Y$ is a smooth divisor in $X$ such that $Z|_Y$ is reduced, then $I_k (D)$ satisfy 
$$I_k (D|_Y) \subseteq I_k(D) \cdot \jepPubScript{O}_Y,$$
with equality when $Y$ is general; see Section~13 for a more general statement.


\noindent
(vii) If $ X \to T$ is a smooth family with a section $s\colon T \to X$, and $D$ is a relative divisor on $X$ that satisfies a suitable condition
(see Section~14 for the precise statement) then 
$$\{ t \in T ~|~ I_k (D_t) \not\subseteq \mathfrak{m}_{s(t)}^q\}$$
is an open subset of $T$, for each $q \geqslant 1$.



\noindent
(viii) If $D_1$ and $D_2$ are $\jepPubBold{Q}$-divisors with supports $Z_1$ and $Z_2$, such that $Z_1 +  Z_2$ is also reduced, then  the subadditivity property
$$I_k (D_1 + D_2) \subseteq I_k (D_1)\cdot I_k (D_2)$$
holds; see Section~15 for a more general statement.
\end{intro-theorem}


For comparison, the list of properties of Hodge ideals in the case when $D$ is reduced is summarized in \cite[\S4]{Popa}. While much of the story carries over to the setting of $\jepPubBold{Q}$-divisors -- besides of course the connection with the classical Hodge theory of the complement $U = X \smallsetminus D$, which only makes sense in the reduced case -- there are a few significant points where the picture becomes more intricate. For instance, the bounds for the generation level of the Hodge filtration can become worse.  
Moreover, we do not know 
whether the inclusions $I_k (D) \subseteq I_{k-1} (D)$ continue to hold for arbitrary $\jepPubBold{Q}$-divisors. New phenomena appear as well: unlike in the case of multiplier ideals, for 
rational numbers $\alpha_1 < \alpha_2$, usually the ideals $I_k (\alpha_1 Z)$ and $I_k (\alpha_2 Z)$ cannot be 
compared for $k \geqslant 1$; see for instance Example 10.5.

It turns out however that most of these issues disappear if one works modulo the ideal of the hypersurface, at least for 
rational multiples of a reduced divisor.
This, as well as other basic facts, is addressed in the sequel \cite{MP3}, which 
studies Hodge ideals from a somewhat different point of view, namely by comparing them to
the (microlocal) $V$-filtration induced on $\jepPubScript{O}_X$ by $h$. This is inspired by the work of Saito \cite{Saito-MLCT} in the reduced case. In the statement below we summarize some of these properties, which complement the results in Theorem \ref{main_list}, but which we do not know how to obtain with the methods of this paper.


\begin{intro-theorem}[\cite{MP3}]\label{other_paper}
Let $D = \alpha Z$, where $Z$ is a reduced divisor and $\alpha \in \jepPubBold{Q}_{> 0}$. Then the following hold:
\begin{enumerate}[label=(\arabic*)]
\item $I_k (D) + \jepPubScript{O}_X(-Z) \subseteq I_{k-1} (D) + \jepPubScript{O}_X(-Z)$ for all $k$.
\item If $\alpha \in (0 , 1]$, then $I_k (D) = \jepPubScript{O}_X \iff k \leqslant \widetilde{\alpha}_Z - \alpha$, where 
$\widetilde{\alpha}_Z$  is the negative of the largest root of the reduced Bernstein-Sato polynomial of $Z$.
\item If $I_{k-1}(D) = \jepPubScript{O}_X$ (we say that $(X, D)$ is $(k-1)$-log canonical),
then $I_{k+1} (D) \subseteq I_k (D)$.
\item Fixing $k$, there exists a finite set of rational numbers $0 = c_0 < c_1 < \cdots < c_s < c_{s+1} = 1$ such that 
for each $0 \leqslant i \leqslant s$ and each $\alpha \in (c_i, c_{i+1}]$ we have 
$$I_k (\alpha Z)\cdot\jepPubScript{O}_Z = I_k (c_{i+1} Z)\cdot\jepPubScript{O}_Z = \text{constant}$$
and such that 
$$I_k (c_{i+1} Z)\cdot\jepPubScript{O}_Z\subsetneq I_k (c_i Z)\cdot\jepPubScript{O}_Z.$$
\end{enumerate}
\end{intro-theorem}

Going back to the description of Hodge ideals by means of log resolutions, the strictness of the Hodge filtration 
for the push-forwards of (summands of) mixed Hodge modules leads to the following local Nakano-type vanishing 
result for $\jepPubBold{Q}$-divisors:

\begin{intro-corollary}\label{intro_local_vanishing}
Let $D$ be an effective $\jepPubBold{Q}$-divisor on a smooth variety $X$ of dimension $n$, and let $f\colon Y\to X$ be a log resolution of $(X,D)$ that is an isomorphism over
$X\smallsetminus {\rm Supp}(D)$. If $E= (f^*D)_{{\rm red}}$, then
$$R^q f_*\big(\Omega_Y^p(\log E)\otimes_{\jepPubScript{O}_Y} \jepPubScript{O}_Y(-\lceil f^*D\rceil)\big)=0\quad\text{for}\quad p+q>n.$$
\end{intro-corollary}

Note that for $p =n$ this is the local vanishing for multiplier ideals \cite[Th. 9.4.1]{Lazarsfeld}, since $E - \lceil f^*D\rceil = - [(1- \varepsilon) f^*D]$ for $0 < \varepsilon \ll 1$. In general, the statement extends the case of reduced $D$ in \cite[Cor. 3]{Saito-LOG} (cf. also \cite[\S{A.5}]{Saito-MLCT}). Unlike \cite[Th. 32.1]{MP1} regarding that case, at the moment we are unable to prove this corollary via more elementary methods.

A different series of applications, given in \cite{MP3}, uses the results proved in this paper together with the relationship between Hodge ideals of $\jepPubBold{Q}$-divisors and the $V$-filtration, in order to describe the behavior of the invariant $\widetilde{\alpha}_Z$ described in Theorem \ref{other_paper} (called the \emph{minimal exponent} of $Z$). For instance, the triviality criterion proved here as Proposition 11.2 
leads to a lower bound \cite[Cor. ~D]{MP3} for  $\widetilde{\alpha}_Z$ in 
terms of invariants on a log resolution, addressing a question of Lichtin and Koll\'{a}r. Moreover, the results
in Theorem \ref{main_list} (vi) and (vii), and Corollary 11.11, lead to effective bounds and to restriction and semicontinuity statements for $\widetilde{\alpha}_Z$, in analogy with well-known properties of log canonical thresholds; for details see \cite[\S6]{MP3}.




